\scanpage{56}
montrerait que la classe de conjugaison des sections détermine la classe de conjugaison de $L_i$, donc $i$.

\uline{Conjecture A} Soit $U$ courbe alg. anabélienne géom. connexe sur corps $K$ de type fini sur $\mathbb{Q}$. Alors toute section de (1) est \deleted{définie} d'un des deux types précédents, i.e. soit définie par un pt de $U(K)$, soit par une section d'une extension (2), avec $i\in I(\bar K)^\Gamma$ i.e. un pt de $\widehat U\setminus U$, rationnel sur $K$.
\margin{Il y a \inserted{\unsure{de plus}} plus d'\illegible{1} les précisions $\pi^{\Gamma^\circ}=(1)$ (premier cas) $\pi^{\Gamma^\circ}=L_i$ (deux. cas)}

Si on admet cette conjecture, alors on en conclurait, par passage à la limite, en considérant le corps des fonctions $L$ de $U$ et $E_L\longrightarrow \unsure{E|B}_K$ ($E_L$ peut être considéré comme un groupe à lacets ``infini'' car une infinité de \inserted{classes de} ss-groupes lacets\inserted{ $L_i$}\ldots) que toute scindage de cette extension provient d'un scindage d'une extension du type (2), avec $i\in I^\Gamma=X(K)$ ($X=\widehat U$). Les classes de tels scindages
\margin{55}
